\section{Platooning}
\label{sec:sota_platooning}

Once the concept of the platoon has been introduced as a cooperative driving application, the state of the art mainly focuses on three issues: how to manage the platoon over time, how to coordinate maneuvers among vehicles, and how to distribute decisions without compromising safety and scalability. The literature shows that the effectiveness of platooning does not depend only on maintaining a short gap, but also on the ability of the whole system to remain consistent during transitions, joins, leaves, and disturbances \cite{segata2014supporting,santini2019platooning,dressler2019cooperative}.

A central aspect is the operational management of the platoon over time. The literature distinguishes between the problem of maintaining the formation in steady-state conditions and the problem of maneuvers, such as platoon formation, \emph{join}, \emph{leave}, \emph{merge} between two platoons, and the handling of exceptional situations. Such maneuvers require explicit coordination protocols and a coherent view of the platoon state; in the presence of packet losses or delays, vehicles may have different local views, with possible consequences for maneuver correctness \cite{segata2014supporting,santini2019platooning}. For this reason, even at the level of high-level maneuvers, communication logic cannot be designed independently of control dynamics.

From an architectural point of view, the proposed approaches include centralized solutions, in which decisions are taken by a coordinating entity, and distributed solutions, in which vehicles negotiate and validate state transitions locally. Distributed approaches are particularly interesting in terms of robustness and scalability, but they require more accurate protocols to guarantee consensus, event ordering, and convergence of decisions \cite{santini2017consensusbased,santini2019platooning}. This issue is also relevant to the thesis, because the \emph{SafeSwitch} mechanism discussed later is conceived as a distributed protocol for degradation management.

Finally, platooning should be framed as a multi-domain cyber-physical system: traffic, control, vehicle dynamics, and communications are tightly coupled. In \cite{segata2016safe}, this interdependence is clearly highlighted both from a theoretical perspective, in terms of control and communication requirements, and from a methodological perspective, through the use of coupled simulators able to represent mobility, networking, and vehicle dynamics together. This approach also forms the methodological foundation of the present work.
